\subsection{The full material-coordinate intertwiner}
\label{x4b:material}
Here is a complete coordinate proof of the material identity used above.
At a fixed permitted real flow time $\tau$, write $X=X_\tau$ on an
open flow domain $\Omega_\tau$, and let $\Omega'_\tau=X(\Omega_\tau)$.
The map
\[
 X^*:C^\infty(\Omega'_\tau)\longrightarrow C^\infty(\Omega_\tau),
 \qquad v\longmapsto v\circ X
\]
is a linear homeomorphism for the usual compact derivative seminorms;
its inverse is composition with the smooth inverse flow. The chain
rule on compact sets proves both continuity assertions. Put
$A_{ij}=\partial_{q_j}X_i$. Differentiating the flow equation gives
$\partial_\tau A=(DU)(X)A$ and $A_0=I$. On every interval of local
invertibility, the determinant derivative is
$\partial_\tau\det A=(\operatorname{div}U)(X)\det A=0$.
The proved zero divergence and the initial value consequently give
$\det A=1$ throughout the flow domain.

For completeness the needed cofactor identity, in the unchanged
index order, is
\[
 (A^{-1})_{ji}=\frac12\epsilon_{iab}\epsilon_{jcd}
                   (\partial_{q_c}X_a)(\partial_{q_d}X_b),\qquad
 \sum_j\partial_{q_j}(A^{-1})_{ji}=0.
\]
The first formula is the cofactor formula with $\det A=1$.
In the derivative of its first factor, the symmetric pair $(j,c)$
contracts against $\epsilon_{jcd}$; in the derivative of its second
factor the symmetric pair $(j,d)$ does so. Each term is zero, proving
the second formula. Therefore every smooth vector field $w$ on
$\Omega'_\tau$ obeys
\[
 \operatorname{div}_q\bigl(A^{-1}(w\circ X)\bigr)
 =\sum_{i,j,k}(A^{-1})_{ji}(\partial_{x_k}w_i)(X)A_{kj}
 =(\operatorname{div}_x w)\circ X.
\]
The gradient chain rule gives
$\nabla_q(v\circ X)=A^{\mathsf T}(\nabla_xv)\circ X$.
With the full matrix $K_\tau=A^{-1}A^{-\mathsf T}$ it follows that
\begin{equation}
 \mathcal L_\tau X^*=X^*\Delta_x,\qquad
 \mathcal L_\tau^2X^*=X^*\Delta_x^2,
 \qquad \mathcal L_\tau=\operatorname{div}_q(K_\tau\nabla_q).
 \label{x4b:operator}
\end{equation}
The second equality is the first one applied to $\Delta_xv$;
all derivatives of $K_\tau$ are part of the composed operator.
This proves the typed isomorphism of the differential operators,
including every spatial coefficient, on the stated smooth domains.

\subsection{The general entire-function residue and the actual fourth coefficient}
\label{x4b:residue}
Let $\mathcal O(\C_s)$ carry its compact-convergence topology. The
original polynomial pullback $\sigma^*$ has domain
$\mathcal O(\C_s)$ and codomain $\mathcal O(\C_q^3)$, and is
injective with the left inverse $\eta^*$ already proved above.
The four coefficients $C_j$ in \eqref{ab:full} are independent of
the particular entire input. The complete chain rule therefore gives
the following equality for every $h\in\mathcal O(\C_s)$:
\begin{align}
 (\Delta_q^2\sigma^*h)(\gamma(z))
    &=\sum_{j=1}^4 C_j(z)h^{(j)}(-z^2/8),\label{x4b:allh}\\
 \mathfrak r(h):=
 \lim_{z\downarrow0}z^2(\Delta_q^2\sigma^*h)(\gamma(z))
    &=36h'(0)+\frac{15}{16}h''(0).
    \label{x4b:residue-map}
\end{align}
To prove the limit without discarding a derivative, all four
$h^{(j)}(-z^2/8)$ converge to $h^{(j)}(0)$; meanwhile
$z^2C_4,z^2C_3\to0$, $z^2C_2\to15/16$ and $z^2C_1\to36$.
These four separate limits prove the displayed equality. This is
a continuous linear map $\mathfrak r:\mathcal O(\C_s)\to\C$:
the Cauchy formula on any fixed circle about zero bounds both
derivative evaluations by its compact supremum.

Its domain can also be restricted to the original $\A$ or the common
space $C=\A\cap\mathcal H$. It is continuous there because weighted
Fourier seminorms bound both derivatives. An explicit right inverse
in those very spaces is
\begin{equation}
 E:\C\longrightarrow C,\qquad
 E(a)(s)=\frac{8a}{15\Xi(0)}s^2\Xi(s),\qquad
 \mathfrak rE(a)=a.
 \label{x4b:section}
\end{equation}
Indeed $s^2\Xi\in C$, $E(a)'(0)=0$ and $E(a)''(0)=16a/15$.
Its scalar-to-function continuity holds for the specified topological
vector spaces. Thus the precise kernel is
$\{h:36h'(0)+(15/16)h''(0)=0\}$, the image is all of $\C$, and
$h\mapsto(h-E\mathfrak r(h),\mathfrak r(h))$ is a continuous
linear isomorphism onto $\ker\mathfrak r\times\C$, with inverse
$(b,a)\mapsto b+E(a)$, on each of these domains.

In particular retain $\mathcal K_t=\mathcal T_{\mathcal D}Z_t$
and the physical incidence factor $a_{ef}$ from \eqref{xat:K}.
Define $\mathcal R(t)=\mathfrak r(\mathcal K_t)$. The exact value is
\begin{align}
 \mathcal R(t)={}&c_0\left(36Z_t'(3)+\frac{15}{16}Z_t''(3)\right)
 -\frac1{336}\left(36Z_t''(3)+\frac{15}{16}Z_t'''(3)\right)\notag\\
 &+c_1\left(36Z_t'(9/2)+\frac{15}{16}Z_t''(9/2)\right)
 +c_2\left(36Z_t'(13/2)+\frac{15}{16}Z_t''(13/2)\right).
 \label{x4b:fullresidue}
\end{align}
Every derivative and shift remains present. The original evenness of
$Z_t$ supplies no deletion of the derivative at any of these nonzero
arguments. Equations \eqref{x4b:allh}--\eqref{x4b:fullresidue} prove
\begin{equation}
 \lim_{z\downarrow0}z^2(\Delta_q^2\Psi^{(4)})(t,\gamma(z))
       =a_{ef}\mathcal R(t).
 \label{x4b:physicalresidue}
\end{equation}
In the material coordinates the exact same number is
\[
 \lim_{\tau\uparrow1/8}(1-8\tau)
 \mathcal L_\tau^2
       [a_{ef}\mathcal K_t(\sigma+\tau)](q_0)=a_{ef}\mathcal R(t),
\]
by \eqref{x4b:operator} and the original inverse
$z=\sqrt{1-8\tau}$, $\tau=(1-z^2)/8$. For $(\nu\Delta_q)^2$
and $(\nu\mathcal L_\tau)^2$ the right sides acquire exactly
$\nu^2$, with the same original $\nu>0$.

This calculation is also a full coordinate morphism through the
invertible fourth-coefficient operator, on all of $\A$:
\[
 \mathfrak r\mathcal T_{\mathcal D}:\A\to\C,
 \qquad
 \ker(\mathfrak r\mathcal T_{\mathcal D})
       =\mathcal T_{\mathcal D}^{-1}(\ker\mathfrak r),\qquad
 (\mathfrak r\mathcal T_{\mathcal D})
       (\mathcal T_{\mathcal D}^{-1}E)=\id_{\C}.
\]
Its coordinate formula is \eqref{x4b:fullresidue} with $Z_t$
replaced by its actual input $h$. The first equality of kernels
is the definition tested in both directions, and the last equality
uses the two proved inverses, so surjectivity and a continuous right
inverse follow without an assertion about pointwise zero preservation.

\subsection{A strictly nonzero transferred residue at negative Newman times}
\label{x4b:negative}
Set
\[
 d_0=c_0+c_1+c_2=\frac{127}{219024}>0,\qquad
 d_1=3c_0-\frac1{336}+\frac92c_1+\frac{13}2c_2=0.
\]
These are the original zeroth and first energy moments of
$\mathcal D$, including its derivative delta. For the unchanged
Fourier multiplier $m(v)$ from Section~\ref{sec:x4-arithmetic},
$m(0)=d_0$ and $m'(0)=id_1=0$. Keep the original kernel $k$,
with
\[
 k(0)=\sum_{n\geq1}\frac\pi2n^2(2\pi n^2-3)e^{-\pi n^2}>0.
\]
The sum converges by Gaussian decay and is strictly positive term
by term. No constant kernel has replaced $k(v)$.

\begin{proposition}[The full transferred asymptotic with a remainder bound]
There is a finite explicitly defined constant $C_*$ below such that,
for every real $T\geq1$,
\begin{equation}
 \mathcal R(-T)=-15d_0k(0)\sqrt\pi\,T^{-3/2}+\mathcal E(T),
 \qquad |\mathcal E(T)|\leq96C_*\sqrt\pi\,T^{-5/2}.
 \label{x4b:asymptotic}
\end{equation}
Consequently $\mathcal R(-T)<0$ for all sufficiently large $T$.
For every original pair with $a_{ef}>0$, the actual transferred
bilaplacian is unbounded along the original escaping trajectory
at each such fixed Newman time.
\end{proposition}
\begin{proof}
Differentiation of the original Fourier integral, justified by its
full weighted bounds, gives
\[
 \mathcal R(-T)=4\int_{\R}e^{-Tv^2/4}k(v)
       \left(-\frac{15}{16}v^2+36iv\right)m(v)\,dv.
\]
The original coefficients are real, so $m(-v)=\overline{m(v)}$.
The kernel $k$ is real and even. Write
\[
 F(v)=k(v)\operatorname{Re}\left[
           \left(-\frac{15}{16}v^2+36iv\right)m(v)\right],
 \qquad c=\frac{15}{16}d_0k(0)>0.
\]
Then $F$ is smooth, real and even. Its Taylor coefficients through
degree two are $F(0)=F'(0)=0$ and $F''(0)/2=-c$.
Indeed the real quadratic coefficient from $36ivm(v)$ is
$-36d_1v^2=0$; the other quadratic coefficient is
$-(15/16)d_0v^2$. This retains and computes the derivative-channel
contribution. Evenness also gives $F'''(0)=0$.

All constants in the following definition are finite:
\[
 K_*:=\sup_{v\in\R}|k(v)|,\quad
 A_*:=|c_0|+c_1+c_2,\quad B_*:=1/336,
\]
\[
 C_*:=\max\left\{
 \frac1{24}\max_{|v|\leq1}|F^{(4)}(v)|,
 \ K_*\left(\frac{15}{16}+36\right)(A_*+B_*)+c
                    \right\}.
\]
The proved two-ended decay of $k$ makes $K_*$ finite, and smoothness
makes the displayed compact derivative maximum finite. For $|v|\leq1$,
Taylor's theorem with its integral remainder gives
$|F(v)+cv^2|\leq C_*v^4$. For $|v|\geq1$, use the complete bound
$|m(v)|\leq A_*+B_*|v|$ and
\[
 |F(v)|\leq K_*\left(\frac{15}{16}v^2+36|v|\right)
                         (A_*+B_*|v|).
\]
Every power on the right is at most $|v|^4$ on that set. The second
entry in the maximum consequently gives the same bound
$|F(v)+cv^2|\leq C_*v^4$ there. It is therefore a global remainder
bound for the original integrand, not only a formal Taylor expansion.
The imaginary integrand is odd and absolutely integrable, so its
integral vanishes by the substitution $v\mapsto-v$.

The Gaussian integral and two derivatives with respect to its positive
parameter give, in the original $v$ coordinate,
\[
 \int_{\R}v^2e^{-Tv^2/4}\,dv=4\sqrt\pi T^{-3/2},\qquad
 \int_{\R}v^4e^{-Tv^2/4}\,dv=24\sqrt\pi T^{-5/2}.
\]
Differentiation is justified on each compact positive parameter
interval by a polynomial times a fixed Gaussian. Integration of the
global remainder bound gives
$\mathcal R(-T)=-16c\sqrt\pi T^{-3/2}+\mathcal E(T)$ with
the stated error. Substitution of $c$ proves \eqref{x4b:asymptotic}.
For example every $T>\max\{1,6C_*/c\}$ makes this value negative.
Equation \eqref{x4b:physicalresidue} then gives unboundedness, since
$z^2>0$ and $a_{ef}\mathcal R(-T)<0$.
\end{proof}

Every such observable is still entire at each finite spatial point.
The exact coordinate map giving the unbounded sequence is
$z\mapsto\gamma(z)=\Gamma(-z/2)$, with inverse $z=1/x$ on its
image and $|\gamma(z)|\geq z^{-1}\to\infty$. In material coordinates
the same sequence is represented by $(\tau,q_0)$ with
$\tau=(1-z^2)/8$, and \eqref{x4b:operator} proves the full operator
correspondence there. The newly proved divergence is thus a property
of these actual scalar observables and their complete diffusion
operators along that image. It supplies no off-critical zero of
$Z_0$ and no finite-point singular classical velocity.
